\documentclass[conference]{IEEEtran}

\usepackage[T1]{fontenc}
\usepackage[utf8]{inputenc}
\usepackage[brazil]{babel}
\usepackage{booktabs}
\usepackage{xurl}
\usepackage[hidelinks]{hyperref}
\usepackage[nopatch=eqnum]{microtype}

\title{Treinamento de um Modelo de Linguagem Autorregressivo em Nível de Caractere com Textos de Machado de Assis}

\author{%
  \IEEEauthorblockN{FABRICIO FERNANDES SANTANA}
  \IEEEauthorblockA{Departamento de Ciência da Computação\\
  Universidade de Brasília\\
  Matrícula: 262114527}
}

\begin{document}
\maketitle

\begin{abstract}
Este trabalho implementa, treina e compara duas configurações de modelos autorregressivos em nível de caractere, treinadas do zero com uma amostra de textos de Machado de Assis. Usando nanoGPT em versão fixada e o mesmo orçamento de atualizações em uma Tesla T4, investigamos como a redução conjunta de camadas, cabeças e dimensão das representações vetoriais se associa à previsão no teste. Em 6.400 janelas amostradas do fluxo de teste, o modelo de referência obteve perda de 1,2397 nats por token e perplexidade 3,455; a configuração reduzida obteve 1,3862 e 4,000. As métricas usam o vocabulário de caracteres, incluindo os tokens especiais de fim de documento e de caractere desconhecido. Incluindo representações posicionais, a configuração reduzida tem 69,8\% menos parâmetros únicos, mas maior perda neste protocolo. O corpus é uma amostra de trabalho, sem alegação de representar a obra completa. Discute-se uma extensão por recuperação de evidências, sem atribuir capacidade conversacional aos modelos atuais.
\end{abstract}

\begin{IEEEkeywords}
modelos de linguagem, Transformer, nanoGPT, Machado de Assis, geração de texto, aprendizado profundo.
\end{IEEEkeywords}

\section{Introdução}

Modelos de linguagem autorregressivos estimam a probabilidade do próximo token, a unidade de representação do texto, a partir de um prefixo. A arquitetura Transformer substitui recorrência por mecanismos de atenção, permitindo relacionar posições distintas da sequência \cite{vaswani2017}. Modelos da família GPT utilizam uma organização apenas com decodificador (\textit{decoder-only}) e máscara causal: a previsão do próximo token depende somente dos tokens já disponíveis \cite{radford2019}. Esse princípio permite estudar modelagem de linguagem com arquiteturas menores que as dos modelos de grande escala.

O objetivo deste trabalho é construir um modelo de linguagem baseado na arquitetura GPT com conteúdo de Machado de Assis, utilizando o nanoGPT como código-base. O trabalho focaliza a implementação de um modelo de referência (\textit{baseline}), a avaliação quantitativa e qualitativa de sua geração e a discussão de uma possível extensão para perguntas sobre o autor. A contribuição é um experimento didático, documentado e auditável em um corpus literário em português, não uma reprodução do GPT-2 original em escala ou tokenização.

O estudo responde à seguinte questão: sob o mesmo corpus, partições e orçamento de 5.000 atualizações, como a redução conjunta de profundidade, número de cabeças e dimensão das representações vetoriais se associa à perda amostral de teste e ao número de parâmetros? Para respondê-la, apresentamos a preparação dos dados, duas configurações de modelo, a comparação quantitativa e a inspeção qualitativa de amostras.

\section{Fundamentação}

O Transformer utiliza atenção para combinar informações de diferentes posições \cite{vaswani2017}. Na variante do nanoGPT, cada token recebe uma representação vetorial aprendida (\textit{embedding}), somada à representação de sua posição. A sequência passa por blocos de atenção causal e redes de propagação direta (\textit{feed-forward}), com normalização e conexões residuais. Após a normalização final, uma projeção produz escores (\textit{logits}) para o vocabulário; a função softmax os transforma nas probabilidades do próximo token \cite{karpathyNanoGPT}.

Para uma cabeça de atenção, a operação é
\begin{equation}
\mathrm{Attention}(Q,K,V)=\mathrm{softmax}\left(\frac{QK^\top}{\sqrt{d_k}}+M\right)V.
\end{equation}
As matrizes $Q$, $K$ e $V$ contêm, respectivamente, consultas, chaves e valores obtidos por projeções aprendidas das representações de entrada. O produto entre consultas e chaves determina os pesos usados para combinar valores; $d_k$ é a dimensão das chaves. A máscara $M$ vale zero nas posições permitidas e $-\infty$ nas futuras, anulando seus pesos após softmax. As cabeças calculam essas relações em paralelo, e suas saídas são concatenadas e projetadas \cite{vaswani2017}. O nanoGPT usa apenas blocos causais, sem o codificador e a atenção cruzada do Transformer original \cite{karpathyNanoGPT}.

O treinamento ajusta os parâmetros $\theta$ minimizando a entropia cruzada. Para $T$ tokens-alvo, com $x_{<t}$ representando o prefixo disponível na janela, a perda média e a perplexidade são
\begin{equation}
L=-\frac{1}{T}\sum_{t=1}^{T}\log p_\theta(x_t\mid x_{<t}),\qquad \mathrm{PPL}=\exp(L).
\end{equation}
O logaritmo natural expressa a perda em nats por token; a perplexidade é adimensional. Neste trabalho, cada token representa um caractere ou um dos símbolos especiais de fim de documento (EOS) e de caractere desconhecido (UNK). Ambos participam da perda quando aparecem nos alvos. Essa convenção é mantida em todas as métricas reportadas.

A tokenização por caracteres segue o exemplo \path{shakespeare_char} do nanoGPT. O tutorial de Karpathy apresenta os componentes passo a passo, enquanto este experimento adapta o código-base para o corpus e seu tokenizador \cite{karpathyNanoGPT,karpathyTutorial}. Uma janela de 256 tokens corresponde a aproximadamente 256 caracteres, descontados os tokens especiais. A tokenização por subpalavras, como BPE, pode representar unidades textuais maiores; suas perplexidades não são diretamente comparáveis às deste estudo, pois muda a unidade de predição \cite{radford2019}.

Brown et al.\ investigam, no GPT-3, aprendizagem em contexto: exemplos e instruções no prefixo orientam a resposta sem atualizar os pesos na aplicação da tarefa \cite{brown2020}. Isso difere do treinamento por gradientes realizado aqui. A escala e o protocolo daquele estudo também diferem dos presentes; prever continuações de texto não demonstra, por si só, aprendizagem em contexto ou capacidade de responder perguntas.

\section{Dados e metodologia}

\subsection{Corpus e partições}

O corpus deriva do pacote \texttt{machado} do NLTK, cujos metadados apontam para a coleção digital do MEC, produzida com participação do NUPILL/UFSC \cite{nltkMachado,mecMachado}. Dos 246 arquivos de texto da fonte, quatro traduções foram excluídas; o mestre de trabalho reúne 242 itens: 112 indexados tratados como autorais e 130 contos avulsos. O inventário inclui contos, crítica, crônicas, romances, teatro e poesia.

Na seleção dos contos avulsos, mediu-se a fração de seus n-gramas distintos de 32 caracteres encontrada em cada uma de sete coletâneas, tomando-se a maior cobertura. A comparação usou o corpo após a nota ``Publicado originalmente'', normalizado em uma sequência alfanumérica sem diacríticos. Coberturas a partir de 0,85 foram classificadas como altas; de 0,40 até abaixo de 0,85, como parciais. Os 130 textos incluídos ficaram abaixo de 0,40 e não tiveram o corpo integral localizado nas coletâneas. Essa triagem não demonstra ineditismo nem ausência de duplicação parcial em outras edições.

O cruzamento com 32 títulos da bibliografia da Academia Brasileira de Letras (ABL) \cite{ablBibliografia} encontrou 24 caminhos diretamente correspondentes, seis casos sem equivalência integral verificada, uma tradução excluída e a coletânea \textit{Correspondência} (1932), não localizada no pacote. Essas contagens descrevem o inventário, não uma equivalência entre arquivos, livros e edições. Assim, os resultados referem-se à amostra de trabalho identificada pelo SHA-256 no manifesto, sem generalização à obra completa.

Os textos da fonte foram decodificados de Windows-1252 e armazenados em UTF-8. A preparação removeu apenas os delimitadores técnicos externos do arquivo mestre, aplicou normalização Unicode NFC, converteu finais de linha para LF, removeu espaços finais e reduziu sequências de mais de duas linhas vazias a duas. Cabeçalhos, notas editoriais, índices e conteúdo literário foram preservados e, portanto, também integram a distribuição aprendida.

A divisão por documento, com semente 20260925, buscou proporções de 80\%/10\%/10\% dos caracteres. Foram alocados 181 documentos ao treino (11.113.713 caracteres), 30 à validação (1.340.673) e 31 ao teste (1.245.563), correspondendo a 81,12\%, 9,79\% e 9,09\%. A diferença decorre dos comprimentos desiguais dos documentos. A verificação de duplicação integral após normalização de espaços não encontrou grupos duplicados; isso não exclui passagens repetidas entre obras distintas.

O tokenizador foi ajustado somente no treino e possui 147 tokens, incluindo EOS e UNK. Houve zero caracteres mapeados para UNK no treino, dois na validação e um no teste. Cada documento é codificado caractere a caractere e terminado por EOS; os fluxos de cada partição são concatenados. Uma janela pode atravessar fronteiras entre documentos, e textos mais longos oferecem mais posições possíveis para amostragem.

\subsection{Modelo e treinamento}

Foi utilizado o nanoGPT, de Andrej Karpathy, sob licença MIT, no commit \texttt{3adf61e154c3} \cite{karpathyNanoGPT}. O baseline tem seis camadas, seis cabeças e dimensão de embedding 384; a condição reduzida tem quatro camadas, quatro cabeças e dimensão 256. Os demais hiperparâmetros são compartilhados, conforme a Tabela~\ref{tab:config}. Ambas as configurações foram treinadas do zero por 5.000 iterações no Colab, com PyTorch 2.11.0+cu128 e Tesla T4.

O treinamento calcula gradientes por retropropagação e atualiza os pesos com AdamW. A semente efetiva do código \texttt{train.py} em uma GPU é 1337; já a preparação dos dados e a geração usam 20260925. A opção \path{always_save_checkpoint=True} grava o estado do treinamento (\textit{checkpoint}) em cada avaliação. Os arquivos \texttt{ckpt.pt} preservados correspondem à iteração 5.000. As curvas registram avaliações a cada 500 iterações e mostram que esse checkpoint teve a menor perda de validação entre os pontos registrados.

Reduzir camadas, cabeças e dimensão simultaneamente constitui uma intervenção composta: os resultados comparam duas configurações, mas não isolam o efeito de cada componente. Há uma execução por condição. O orçamento iguala atualizações e tokens processados, não operações de ponto flutuante (FLOPs) nem tempo de execução.

\begin{table}[t]
\caption{Hiperparâmetros do treinamento.}
\label{tab:config}
\centering
\begin{tabular}{ll}
\toprule
Item & Configuração \\
\midrule
Inicialização & do zero \\
Camadas/cabeças/dimensão & 6/6/384; reduzido 4/4/256 \\
Contexto / lote & 256 tokens / 64 \\
Dropout / bias & 0,1 / desativado \\
Otimizador & AdamW; $\beta_1=0{,}9$, $\beta_2=0{,}95$ \\
Taxa inicial / mínima & $3\cdot10^{-4}$ / $3\cdot10^{-5}$ \\
Aquecimento / decaimento & 100 / 5.000 iterações \\
Weight decay / clipping & 0,1 / 1,0 \\
Precisão / compilação & float16 / desativada \\
Iterações / GPU & 5.000 / Tesla T4 \\
\bottomrule
\end{tabular}
\end{table}

\subsection{Avaliação}

A avaliação abre os documentos de teste somente após o treinamento de cada modelo. Usa 200 lotes de 32 janelas de 256 tokens, sorteadas do fluxo concatenado de teste com semente NumPy 20260926. Cada lote contém entradas e alvos deslocados em um token. Calcula-se a média das 200 perdas de lote e aplica-se a exponencial para obter a perplexidade. Trata-se de uma estimativa amostral, não de uma varredura exaustiva nem de uma média por obra. Os inícios podem se repetir ou gerar janelas sobrepostas; as 6.400 sequências não são necessariamente independentes. A semente torna o sorteio reproduzível, mas não remove a incerteza amostral.

Ambas as condições foram avaliadas no Colab em Tesla T4, com PyTorch 2.11.0+cu128, a mesma semente de teste e o mesmo protocolo. O conjunto de teste contém uma ocorrência do caractere ``½'', ausente do vocabulário ajustado no treino; ela foi mapeada para UNK e permaneceu como alvo na avaliação. O resultado é, portanto, uma comparação pareada no ambiente, embora baseada em uma amostra de janelas.

A geração usa o prefixo ``Era '' (incluindo o espaço final), temperatura 0,8, top-k 40 e semente 20260925. São gerados 700 novos tokens, e o texto apresentado é truncado no primeiro EOS, quando presente. A temperatura reescala os escores antes da amostragem; top-k restringe a distribuição aos 40 candidatos de maior pontuação. O prefixo é um início narrativo, não uma pergunta. As amostras servem à inspeção qualitativa e não substituem avaliação humana sistemática de qualidade, autoria ou factualidade.

\section{Resultados}

\begin{table}[t]
\caption{Avaliação amostral. Perda em nats por token do vocabulário de caracteres, incluindo EOS e UNK; PPL calculada sobre os mesmos tokens. Parâmetros incluem posições.}
\label{tab:comparison}
\centering
\begin{tabular}{lrrr}
\toprule
Modelo & Parâmetros totais & Perda & PPL \\
\midrule
Baseline (6/6/384) & 10,78 M & 1,2397 & 3,455 \\
Reduzido (4/4/256) & 3,25 M & 1,3862 & 4,000 \\
\bottomrule
\end{tabular}
\end{table}

Os dois checkpoints foram carregados com sucesso e registraram 5.000 iterações. A contagem total inclui embeddings posicionais e conta uma vez os pesos dos embeddings de token compartilhados com a projeção de saída: 10.776.576 no baseline e 3.251.200 na variante, redução de 69,8\%. Sem posições, o nanoGPT informa 10.678.272 e 3.185.664 parâmetros. A perda amostral aumentou 0,1465 nats por token (11,8\%) na configuração menor, e a perplexidade aumentou 0,545 (15,8\%). As diferenças foram calculadas com a precisão integral dos resultados armazenados, antes do arredondamento da Tabela~\ref{tab:comparison}.

Na amostra baseline, “Era uma coisa de Fernando recebida na sua imaginação” combina palavras e pontuação plausíveis, mas não estabelece relação semântica clara; a continuação mistura datas e nomes. Na amostra reduzida, “as duas amizas” e “a cadeira um ar de tendência” ilustram grafia inventada e continuidade instável. As duas gerações reproduzem pontuação e moldes de diálogo reconhecíveis, sem manter coerência narrativa. Esta inspeção é ilustrativa, não uma avaliação humana sistemática, e não sustenta uma conclusão global sobre qualidade literária.

\subsection{Comparação experimental}

Sob o protocolo compartilhado, o baseline apresenta menor perda e perplexidade amostrais; a condição reduzida utiliza menos parâmetros. A relação exponencial entre as métricas permite expressar a diferença de perdas como uma razão de perplexidades:
\begin{equation}
\frac{\mathrm{PPL}_{\mathrm{red}}}{\mathrm{PPL}_{\mathrm{base}}}
=\exp(L_{\mathrm{red}}-L_{\mathrm{base}})\approx 1{,}158.
\end{equation}
Portanto, a perplexidade amostral da variante é cerca de 15,8\% maior. Sem erro-padrão nem repetições com sementes independentes, a incerteza não está quantificada. Os valores descrevem estas execuções, mas não estimam com precisão o desempenho esperado em novos treinamentos.

A redução do número de pesos implica menor armazenamento dos parâmetros sob a mesma precisão numérica. Não foram medidos latência, energia, taxa de processamento ou pico de memória em protocolo controlado. Por isso, a redução paramétrica de 69,8\% não deve ser interpretada como ganho equivalente de velocidade ou economia computacional medida.

\section{Discussão}

Os resultados indicam um compromisso entre número de parâmetros e qualidade preditiva: o modelo menor requer menos pesos, mas atribui menor probabilidade média geométrica aos tokens-alvo do teste. A menor perda do baseline não se traduz, nas amostras apresentadas, em continuidade narrativa consistente. Assim, a métrica de previsão local e a inspeção da geração descrevem aspectos distintos do comportamento dos modelos.

O corpus e a divisão são tão importantes quanto a arquitetura. Duplicatas ou sobreposição entre documentos podem reduzir artificialmente a dificuldade do teste; por isso, a auditoria de proveniência e a divisão por obra devem permanecer parte do protocolo. Estudos sobre deduplicação mostram que sobreposição de dados pode afetar tanto a avaliação quanto a memorização de modelos de linguagem \cite{lee2022}.

As curvas de treino e validação das execuções atuais estão preservadas junto aos logs no Drive; a perda de validação diminuiu nos pontos registrados até a iteração 5.000. Esse resultado permite descrever a trajetória amostrada, embora não substitua repetições independentes para estimar variabilidade entre treinamentos. O checkpoint reportado é o da iteração 5.000, que também foi o de menor perda entre os pontos de avaliação registrados.

Há três limites principais à interpretação. Primeiro, a composição do pacote NLTK e a seleção heurística dos contos delimitam a representatividade da amostra. Segundo, a estimativa por janelas sobrepostas não quantifica a incerteza da perda; além disso, a ocorrência de ``½'' foi representada pelo token UNK. Terceiro, a intervenção arquitetural composta e a execução única por condição impedem separar seus componentes da variabilidade entre treinamentos. As duas condições atuais foram avaliadas no mesmo ambiente e protocolo; resultados de rodadas históricas em outros ambientes não são usados como comparação principal.

Há também uma limitação temporal: o resultado do baseline já era conhecido quando se decidiu executar a variante. O teste não foi usado para calcular gradientes, mas a escolha da segunda configuração não foi cega aos resultados anteriores. A comparação deve, portanto, ser interpretada como exploratória. Em novos estudos, as configurações podem ser fixadas antes de consultar o teste, com decisões de arquitetura baseadas na validação.

\section{Extensão para perguntas sobre Machado de Assis}

Os modelos autorregressivos treinados apenas para continuação de texto não devem ser apresentados como sistemas de perguntas e respostas. Uma extensão possível seria combinar recuperação de passagens e geração. As obras seriam segmentadas e indexadas com título, capítulo e posição; diante de uma pergunta, o sistema recuperaria trechos relevantes e os forneceria ao gerador junto com instruções para responder com base nas evidências. A resposta deveria citar os trechos recuperados ou declarar insuficiência de evidência.

Por exemplo, diante de ``Como Bentinho interpreta a ausência de Capitu?'', o sistema deveria localizar passagens pertinentes antes de formular a resposta. O contexto atual de 256 tokens de caracteres é restrito para acomodar pergunta, instrução e vários excertos. Seria necessário ampliar o contexto ou substituir/adaptar o gerador, além de avaliar sua capacidade de usar as evidências. Essa proposta se relaciona à geração aumentada por recuperação (RAG), mas recuperar trechos e inseri-los no prefixo não reproduz o treinamento conjunto do sistema original de Lewis et al.\ \cite{lewis2020rag}.

Um primeiro experimento dessa extensão poderia usar respostas extrativas, compostas pelas próprias passagens recuperadas, como referência antes de introduzir geração. A avaliação mediria precisão e cobertura dos trechos relevantes, apoio textual das respostas, atendimento à pergunta e abstenção apropriada, com perguntas anotadas e critérios definidos antes do teste. Esta é uma proposta, não uma funcionalidade implementada.

\section{Conclusão}

Foram implementadas e comparadas duas configurações autorregressivas em nível de caractere. Sob o mesmo orçamento de atualizações, a variante reduzida tem 69,8\% menos parâmetros únicos, incluindo posições, mas apresenta perda amostral 11,8\% maior e perplexidade 15,8\% maior. A comparação é descritiva e restrita às execuções e à amostra selecionada. As gerações exibem padrões locais de linguagem e falhas de coerência, sem demonstrar capacidade conversacional.

O experimento articula preparação rastreável de um corpus literário, treinamento do zero e avaliação dos limites preditivos e qualitativos de modelos pequenos. A preservação de curvas de validação e a avaliação com repetições fortaleceriam estudos futuros. Para perguntas sobre Machado de Assis, uma extensão possível é combinar recuperação de trechos e geração com evidências explícitas, avaliando separadamente os dois componentes.

\section*{Reprodutibilidade e atribuição}

Os arquivos \texttt{model.py}, \texttt{train.py} e utilitários são do nanoGPT, de Andrej Karpathy, sob licença MIT; a configuração de caracteres deriva de \texttt{shakespeare\_char} \cite{karpathyNanoGPT}. Os notebooks do projeto acrescentam auditoria e preparação do corpus, tokenizador, partições, configurações, avaliação e exportação. Não se atribui autoria própria ao código-base.

Os notebooks executados e as instruções de reprodução estão no \href{https://github.com/fabriciosantana/unb/blob/main/01-ppgi0034-redes-neurais-e-aprendizado-profundo/projetos/machado-assis/README.md}{README do projeto}. Os quatro notebooks abrem no Colab: \href{https://colab.research.google.com/github/fabriciosantana/unb/blob/main/01-ppgi0034-redes-neurais-e-aprendizado-profundo/projetos/machado-assis/notebooks/01_construcao_auditoria_corpus.ipynb}{01 auditoria}, \href{https://colab.research.google.com/github/fabriciosantana/unb/blob/main/01-ppgi0034-redes-neurais-e-aprendizado-profundo/projetos/machado-assis/notebooks/02_preparacao_dados_modelagem.ipynb}{02 preparação}, \href{https://colab.research.google.com/github/fabriciosantana/unb/blob/main/01-ppgi0034-redes-neurais-e-aprendizado-profundo/projetos/machado-assis/notebooks/03_baseline_gpt_caractere_nanogpt.ipynb}{03 baseline} e \href{https://colab.research.google.com/github/fabriciosantana/unb/blob/main/01-ppgi0034-redes-neurais-e-aprendizado-profundo/projetos/machado-assis/notebooks/04_comparativo_gpt_caractere_reduzido.ipynb}{04 comparativo}. Dados e resultados estão na pasta compartilhada do \href{https://drive.google.com/drive/folders/1VgfGCoNn-FatGc0q5tSDLwDOgLNgdkGa}{Google Drive}; o \href{https://drive.google.com/file/d/1F5IyjmACoBWa8HJQLv8BQb8EuwnpjMyx/view?usp=drivesdk}{PDF final} também está disponível nela.

\noindent Página de entrada do projeto: \href{https://github.com/fabriciosantana/unb/blob/main/01-ppgi0034-redes-neurais-e-aprendizado-profundo/projetos/machado-assis/README.md}{README e instruções de avaliação}.

\bibliographystyle{IEEEtran}
\bibliography{referências/bibliografia}

\end{document}
